\section{Linearly many pairs with logarithmic neighborhood discrepancy}
\label{geo:section}

For two distinct vertices $u,v$ of a finite simple graph $G$, let
\[
 d_G(u,v)=\bigl|\{w\notin\{u,v\}:\text{$w$ is adjacent to exactly one of $u,v$}\}\bigr|.
\]
Thus $d_G(u,v)$ counts the entries on which their adjacency rows disagree
after the two endpoint positions have been removed. It is not a graph
distance. Write $H_D(G)$ for the graph on $V(G)$ in which $uv$ is an
edge if and only if $d_G(u,v)\le D$; this says nothing about whether
$uv$ is an edge of $G$. The ordinary deck is the multiset of unlabelled
graphs $G-v$, with one occurrence for each vertex. Nonreconstructible
means that some nonisomorphic graph has the same ordinary deck.

\begin{theorem}[A linear matching of near-twins]\label{geo:matching}
Let $G$ be a nonreconstructible graph of order $n\ge 2^{20}$ and put
$D=\lceil320\log_2 n\rceil$. Every vertex cover of $H_D(G)$ has more
than $n/160$ vertices. In particular, every inclusion-maximal matching
of $H_D(G)$ has more than $n/320$ edges.
\end{theorem}

The conclusion concerns distinct \emph{endpoints} in one parent. It does
not say that the sets of vertices distinguishing the matched pairs are
disjoint, or that a surviving mate can be identified on an individual
card. The proof proceeds by contraposition. A small cover leaves a large
set of mutually distinct exterior rows. Sampling exposes enough of their
differences to obtain landmarks. The difficulty is to count all placements
of those landmarks without paying for every permutation of the cover.

\subsection{A landmark criterion requiring no distinguished copy}

A named graph here is a graph on a fixed set of positions. For such a
graph $F$, let $c_X(F)$ be the number of injections of its positions
into $X$ preserving both edges and nonedges. If $|F|<n=|X|$, then
\begin{equation}
 c_X(F)=\frac{1}{n-|F|}\sum_{v\in V(X)}c_{X-v}(F).
 \label{geo:proper}
\end{equation}
Indeed, an injection survives precisely the deletions outside its image.
Thus comparing named pattern counts requires no names on an input card.

Fix an ordered tuple $L$ of $r<n$ distinct vertices of $G$. An exterior
vertex has a binary \emph{word} giving its adjacency to $L$ in order.
Call $L$ resolving if these words are distinct among vertices outside
$L$. No assumption is made about other induced copies of $G[L]$.

\begin{lemma}[Resolving landmarks]\label{geo:landmarks}
Let $L$ be resolving, let $F=G[L]$ be named by the tuple, and let
$N=n-r$. If
\begin{equation}
 \frac{c_G(F)}{|\operatorname{Aut}(G)|}<2^{N/2-1},
 \label{geo:criterion}
\end{equation}
then the ordinary deck determines $G$ up to isomorphism.
\end{lemma}

\begin{proof}
In an arbitrary host $X$ of order $n$, fix one induced injection
$j:F\to X$. Read the exterior words relative to $j$. Discard positions
with a word absent from the target $G$. Every remaining position has a
unique target vertex as its name, though several positions in $X$ may
receive the same name. Define a conflict graph $C_j$ on these positions:
two conflict if their names coincide, or if the names are distinct and
their adjacency in $G$ disagrees with the adjacency of the positions
in $X$. An independent set therefore gives a set $S$ of distinct target
names and exactly one induced injection of $G[L\cup S]$ extending $j$.
Conversely, each such injection gives one independent set in the
conflict graph of its own root placement.

Put $I(C;z)=\sum_{A\text{ independent in }C}z^{|A|}$ and
$P_X(z)=\sum_j I(C_j;z)$, summing over \emph{all} induced placements
of $F$. The bijection just described gives
\begin{equation}
 [z^k]P_X(z)=
 \sum_{\substack{S\subseteq V(G)\setminus L\\|S|=k}}c_X(G[L\cup S]).
 \label{geo:coeff}
\end{equation}
By \eqref{geo:proper}, the coefficients below degree $N$ agree between
hosts with the same deck. At degree $N$ the only possible embeddings are
isomorphisms from the full target. Thus a nonisomorphic host $X$ with
the same deck as $G$ would satisfy
\[
 P_G(z)-P_X(z)=|\operatorname{Aut}(G)|z^N.
\]

Set $b=1+\sqrt2$, $a=2+\sqrt2$. For every graph $C$ on $m$ vertices,
$|I(C;-a)|\le b^m$. For a union of isolated vertices and isolated edges,
use multiplicativity and the factors $1-a=-b$ and $1-2a=-b^2$. Otherwise
split independent sets according to inclusion of a vertex of degree
$d\ge2$. Induction bounds the result by $b^{m-1}+a b^{m-d-1}\le b^m$,
since $b^2+a\le b^3$. This includes the empty graph as the base case.
Consequently $|P_X(-a)|\le c_X(F)b^N$ and likewise for $G$. Since
$c_X(F)=c_G(F)$, the displayed polynomial difference would force
\[
 \frac{c_G(F)}{|\operatorname{Aut}(G)|}
 \ge \frac12(a/b)^N=2^{N/2-1},
\]
contrary to \eqref{geo:criterion}. The tuple $L$ was used only in the
comparison proof; its location was not provided in any card.
\end{proof}

\subsection{Keeping the cover anonymous}

We use one external group-theoretic input. A permutation's support is
the set of points it moves. By Kempe, Pyber and Shalev
\cite[Theorem~A(2)]{KPS2007}, a permutation group on $t\ge2$ points
in which every nonidentity element has support of size at least
$\log_2 t$ has order at most $2^{10t}$. Transitivity is not required.

Let $U\subseteq V(G)$ and $B=V(G)\setminus U$. Partition $U$ by equality
of its full adjacency words on $B$. There is a set $W_0\subseteq B$ of
size at most $|U|$ that defines the same partition. Start with no test
coordinates. Whenever two unequal full words remain together, add a
coordinate on which they differ. Each addition increases the number of
classes. For $U=\varnothing$ take $W_0=\varnothing$.

\begin{lemma}[An anonymous-cover automorphism bound]\label{geo:virtual}
Let $T\supseteq U\cup W_0$, $t=|T|\ge2$, and $B_T=T\setminus U$.
Assume every distinct $b,b'\in B_T$ have at least $\log_2 t$
distinguishers in $G[T]$. Then
\begin{equation}
 \frac{|\operatorname{Aut}(G[T])|}{|\operatorname{Aut}(G)|}
 \le 3^t 2^{10t}.
 \label{geo:quotient}
\end{equation}
\end{lemma}

\begin{proof}
The subgroup $A$ preserving $U$ and $W_0$ separately as sets has index
at most $3^t$: count the possible images of the three-part coloring of
$T$. Its elements permute $B_T$ and permute the classes of $U$, since
equality of words on the setwise preserved $W_0$ detects the full
partition. For each class $C$ create $|C|$ formal positions
$(C,1),\dots,(C,|C|)$. Act on $B_T$ as on actual vertices and on these
new positions by $(C,i)\mapsto(gC,i)$. This is an action on exactly $t$
points. It does not assume that $g$ respects any numbering of the
actual vertices inside $C$.

Its kernel consists precisely of elements fixing $B_T$ pointwise: such
an element fixes $W_0$ pointwise and hence preserves every class. Each
kernel element extends, by the identity off $T$, to an automorphism of
$G$. In fact it fixes the $B_T$ side, and vertices exchanged inside a
class have equal incidences to every vertex of $B$, including those
outside $T$. These extensions are distinct. The kernel therefore has
order at most $|\operatorname{Aut}(G)|$.

A nonidentity element of the image moves some actual $b\in B_T$ to
$b'\ne b$. Any distinguisher in $B_T$ must move. If a vertex of a class
$C$ distinguishes $b,b'$, then all vertices of $C$ do, and $gC\ne C$:
otherwise invariance of the incidences to the class would give the same
bit at $b$ and $b'$. All $|C|$ formal positions of such a class move.
Hence the image element moves at least $d_{G[T]}(b,b')$ points. Apply
the quoted motion bound to the image (the trivial image also satisfies
the order bound), then multiply by the kernel and index estimates.
\end{proof}

\subsection{Sampling the landmarks and closing the count}

\begin{proof}[Proof of Theorem~\ref{geo:matching}]
Suppose $U$ covers $H_D(G)$ and $u=|U|\le n/160$. Put
$\ell=\log_2 n$. Distinct vertices of $B=V(G)\setminus U$ have more
than $D$ distinguishers in $G$, which may include vertices of $U$.
Sample a set $Z$ by including \emph{every} vertex independently with
probability $1/80$. For each fixed pair in $B$, its number $X$ of sampled
distinguishers is binomial with mean $\mu\ge4\ell$. Markov's inequality,
applied to $4^{-X}$, gives
\[
 \Pr(X<\ell)\le 4^\ell\mathbb E[4^{-X}]
 \le \exp\bigl(\ell\ln4-3\mu/4\bigr)
 \le \exp\bigl(-(3-\ln4)\ell\bigr).
\]
A union bound over fewer than $n^2$ pairs is at most $n^{-0.328}<1/26$
for $n\ge2^{20}$, since $(3-\ln4)/\ln2-2>0.328$.
Also $\Pr(|Z|>13n/1000)\le25/26$ by Markov's inequality. Thus there is
one sample for which both bounds hold; independence of these success
events was not used.

Use $W_0$ as above and take $T=U\cup W_0\cup Z$. Then
$t=|T|\le2u+13n/1000\le0.0255n<0.026n$. Every two exterior vertices
are distinguished by $T$, as are all distinct vertices of $B_T$ by at
least $\ell$ vertices of $T$. Thus the exterior words relative to $T$
are distinct and Lemma~\ref{geo:virtual} applies, since $\ell\ge\log_2 t$.
The automorphism bound is immediate also for $t=0,1$.

Temporarily order $T$ and put $F=G[T]$ with these names. Each choice of
its image set supports either no induced injection or precisely
$|\operatorname{Aut}(G[T])|$ of them. By \eqref{geo:quotient},
\[
 \frac{c_G(F)}{|\operatorname{Aut}(G)|}
 \le \binom nt\,3^t 2^{10t}.
\]
With $H_2(x)=-x\log_2 x-(1-x)\log_2(1-x)$ (and the continuous endpoint
values), $\binom nt\le2^{n H_2(t/n)}$ by the binomial theorem. The
base-two logarithm of the right-hand side is at most
\[
 n\bigl(H_2(0.026)+0.026(10+\log_2 3)\bigr)<0.48n
 <(n-t)/2-1.
\]
For the last inequality use $(n-t)/2-1\ge0.487n-1$ in the asserted
range. The landmark criterion \eqref{geo:criterion} now holds with
$N=n-t$, so $G$ is reconstructible. This rules out every such
cover in a counterexample. The endpoints of any inclusion-maximal
matching form a vertex cover: an uncovered edge could be added to the
matching. Dividing the strict cover bound by two proves the claim.
\end{proof}
